These relations being isobaric, the quotient is a graded @@M0@@-algebra @@M1@@. We define a morphism @@M2@@ by:
@@M3@@
The imposed relations show that this is a Chern theory for @@M4@@ with values in @@M5@@. If in addition we have a Chern theory @@M6@@, with values in a graded @@M7@@-algebra @@M8@@, there is a unique homomorphism: @@M9@@, such that @@M10@@ factors through @@M11@@; it is defined by sending, for every @@M12@@ and every @@M13@@, the element @@M14@@ to the @@M15@@-th coefficient of @@M16@@ (which is well defined since @@M17@@ is a @@M18@@-homomorphism).

Now suppose that @@M19@@ is of characteristic zero and @@M20@@ has discrete filtration; let @@M21@@ be such that @@M22@@. Then the graded algebra @@M23@@ is zero in degrees @@M24@@. To prove this, we must show that for @@M25@@ and @@M26@@, we have:
@@M27@@
Now, by hypothesis, we have:
@@M28@@
Applying the @@M29@@-homomorphism @@M30@@, we obtain:
@@M31@@
Now, by (6.6.4) and (6.5.3), this implies:
@@M32@@
Since @@M33@@ is of characteristic zero, we have the desired result.

The universal property of @@M34@@, applied to the canonical homomorphism @@M35@@, gives a commutative diagram:
